\chapter{Приближение якобиана}\label{ch:jacobian}
\section{Постановка задачи}
%% ============================================================
 
Рассмотрим задачу решения нелинейной системы
\begin{equation}\label{eq:system}
  F(x) = 0, \qquad F\colon \mathbb{R}^n \to \mathbb{R}^n.
\end{equation}
Квазиньютоновские методы строят последовательность $\{x_k\}$ по правилу
\begin{equation}\label{eq:step}
  x_{k+1} = x_k + d_k, \qquad B_k\,d_k = -F(x_k),
\end{equation}
где $B_k \in \mathbb{R}^{n \times n}$~--- аппроксимация якобиана $J(x_k) = F'(x_k)$.
 
Обозначим
\[
  s_k = x_{k+1} - x_k = d_k, \qquad y_k = F(x_{k+1}) - F(x_k).
\]
 
%% ============================================================
\section{Общая формула: \texorpdfstring{$m$}{m} секущих условий}
%% ============================================================

Пусть заданы $m$ секущих условий: $B_{k+1} S_m = Y_m$, где
\[
  S_m = [s_{k-m+1}\mid\cdots\mid s_k] \in \mathbb{R}^{n \times m}, \qquad
  Y_m = [y_{k-m+1}\mid\cdots\mid y_k] \in \mathbb{R}^{n \times m}.
\]
\emph{Соглашение об обозначениях.} Параметр $m$ используется только в
настоящей главе как ранг общей мультисекантной формулы; начиная
с~гл.~\ref{ch:sp-broyden-theory} обозначение меняется на $p$~--- глубину
сохраняемых \emph{прошлых} секущих условий, при этом число столбцов
окна равно $p{+}1$ за счёт текущей секущей $s_k$
($S_p\in\RR^{n\times(p+1)}$); см.~подробнее раздел <<Обозначения и сокращения>>.

Семейство обновлений, удовлетворяющих всем $m$ секущим условиям:
\begin{equation}\label{eq:general}
  \boxed{
    B_{k+1} = \widehat{B}_k + (Y_m - \widehat{B}_k S_m)\,(V_m^\top S_m)^{-1}\,V_m^\top
  }
\end{equation}
где $\widehat{B}_k \in \mathbb{R}^{n \times n}$~--- базовая матрица,
$V_m \in \mathbb{R}^{n \times m}$~--- матрица направлений ($V_m^\top S_m$ невырождена).

\textbf{Проверка:}
$B_{k+1} S_m = \widehat{B}_k S_m + (Y_m - \widehat{B}_k S_m)(V_m^\top S_m)^{-1}V_m^\top S_m
= \widehat{B}_k S_m + (Y_m - \widehat{B}_k S_m) = Y_m$. \quad
Условия выполнены при \emph{любом} допустимом $V_m$.

Обновление $\Delta_k = B_{k+1} - \widehat{B}_k$ имеет ранг $\leq m$ (произведение
$n \times m$ и $m \times n$ матриц).


%% ============================================================
\section{Частный случай: \texorpdfstring{$m = 1$}{m = 1}}
%% ============================================================

При $m = 1$ имеем $S_m = s_k$, $Y_m = y_k$, $V_m = v_k \in \mathbb{R}^n$,
и формула~\eqref{eq:general} принимает вид:
\begin{equation}\label{eq:rank1}
  B_{k+1} = \widehat{B}_k + \frac{(y_k - \widehat{B}_k s_k)\,v_k^\top}{v_k^\top s_k}.
\end{equation}
Это ранг-1 обновление, удовлетворяющее единственному секущему условию
$B_{k+1} s_k = y_k$.

При дополнительном ограничении $v_k^\top s_k = 1$ знаменатель исчезает:
\begin{equation}\label{eq:rank1_normalized}
  B_{k+1} = \widehat{B}_k + (y_k - \widehat{B}_k s_k)\,v_k^\top.
\end{equation}


%% ============================================================
\section{Разрушение секущих условий при ранг-1 обновлении}
%% ============================================================

Прежде чем переходить к выбору матрицы направлений $V_m$ в общей
формуле~\eqref{eq:general}, отметим характерную патологию ранг-1
ветви~\eqref{eq:rank1}: текущее секущее условие $B_{k+1}s_k = y_k$
выполняется по построению, однако \emph{прошлые} секущие в общем случае
нарушаются. Это и есть основная мотивация перехода от классической
ранг-1 формулы к мультисекантным расширениям~\cite{schnabel1983},
рассматриваемым в гл.~\ref{ch:sp-broyden-theory}.

\begin{proposition}\label{prop:destruction}
Пусть $B_{k+1}$ задано ранг-1 обновлением~\eqref{eq:rank1}
с произвольным $v_k$ ($v_k^\top s_k \neq 0$).
Тогда для любого $j < k$:
\[
  B_{k+1} s_j = B_k s_j + (y_k - B_k s_k)\,\frac{v_k^\top s_j}{v_k^\top s_k}.
\]
В частности, $B_{k+1} s_j = B_k s_j$ тогда и только тогда, когда $v_k^\top s_j = 0$.
\end{proposition}

\begin{proof}
Прямое вычисление:
\[
  B_{k+1} s_j
  = \biggl(B_k + \frac{(y_k - B_k s_k)\,v_k^\top}{v_k^\top s_k}\biggr) s_j
  = B_k s_j + (y_k - B_k s_k)\,\frac{v_k^\top s_j}{v_k^\top s_k}.
\]
Второе слагаемое зануляется тогда и только тогда, когда $v_k^\top s_j = 0$.
\end{proof}

\begin{corollary}\label{cor:broyden_destruction}
В классическом методе Бройдена ($v_k = s_k$), если на предыдущем шаге
выполнялось $B_k s_j = y_j$, то $B_{k+1} s_j = y_j$ тогда и только тогда,
когда $s_k^\top s_j = 0$. В общем случае равенство нарушается.
\end{corollary}


%% ============================================================
\section{Роль \texorpdfstring{$V_m$}{V\_m}: что минимизирует обновление}
%% ============================================================

Выбор $V_m$ при фиксированных $\widehat{B}_k$, $S_m$, $Y_m$ определяет
\emph{форму} обновления $\Delta_k$. Различные $V_m$ удовлетворяют
одним и тем же секущим условиям, но отличаются по норме изменения
$\|\Delta_k\|_F = \|B_{k+1} - \widehat{B}_k\|_F$.

\subsection{\texorpdfstring{$V_m = S_m$}{V\_m = S\_m}: минимум нормы Фробениуса при \texorpdfstring{$m$}{m} секущих условиях}

Обновление принимает вид:
\begin{equation}\label{eq:Vm_Sm}
  B_{k+1} = \widehat{B}_k + (Y_m - \widehat{B}_k S_m)\,(S_m^\top S_m)^{-1}\,S_m^\top.
\end{equation}
Это \textbf{единственное} решение задачи
$\min_{B}\|B - \widehat{B}_k\|_F$ при $B\,S_m = Y_m$
(проекция $\widehat{B}_k$ на аффинное подпространство, задаваемое секущими
условиями, по норме Фробениуса; строгая выпуклость гарантирует единственность).

Ранг обновления $\leq m$. При большом $m$ это может быть существенно:
обновление ранга $m$ требует $O(nm)$ операций на хранение и $O(n^2 m)$ на
применение, что при $m \gg 1$ дороже ранг-1 обновления.

\subsection{\texorpdfstring{$m = 1$, $v_k = s_k$}{m = 1, v\_k = s\_k}: минимум нормы Фробениуса при 1 секущем условии}

Частный случай предыдущего при $m = 1$:
\begin{equation}\label{eq:good_broyden}
  B_{k+1} = \widehat{B}_k + \frac{(y_k - \widehat{B}_k s_k)\,s_k^\top}{\|s_k\|^2}.
\end{equation}
Это решение $\min_B \|B - \widehat{B}_k\|_F$ при $B\,s_k = y_k$. При
$\widehat{B}_k = B_k$~--- классический \textbf{метод Бройдена}
(``good Broyden''~\cite{broyden1965}); ранг~1, затраты $O(n^2)$ на шаг.

\subsection{\texorpdfstring{$m = 1$, $v_k = y_k$}{m = 1, v\_k = y\_k}: другое ранг-1 обновление}

\begin{equation}\label{eq:bad_broyden_direct}
  B_{k+1} = \widehat{B}_k + \frac{(y_k - \widehat{B}_k s_k)\,y_k^\top}{y_k^\top s_k}.
\end{equation}
Удовлетворяет секущему условию, но \textbf{не минимизирует}
$\|B_{k+1} - \widehat{B}_k\|_F$. Связано со \textbf{вторым методом Бройдена}
(``bad Broyden''~\cite{broyden1965}): формула~\eqref{eq:bad_broyden_direct} для прямого
якобиана $B_k$ соответствует минимуму $\|H_{k+1} - H_k\|_F$ при
$H_{k+1} y_k = s_k$ для обратного якобиана $H_k \approx J^{-1}$~\cite{dennis1996}.

\subsection{\texorpdfstring{$m = 1$, произвольный $v_k$}{m = 1, произвольный v\_k}: общий ранг-1 случай}

Любой вектор $v_k$ с $v_k^\top s_k \neq 0$ даёт допустимое ранг-1 обновление.
Gay~\cite{gay1979} показал, что для \emph{линейных} систем выбор $v_k$ не влияет на
итоговую сходимость (метод заканчивается за $2n$ шагов при любом $v_k$).
Для нелинейных систем различные $v_k$ дают различное поведение;
на практике используется почти исключительно $v_k = s_k$ (``good Broyden'').


%% ============================================================
\section{Роль базовой матрицы \texorpdfstring{$\widehat{B}_k$}{B\_k с шапкой}}
%% ============================================================

\subsection{\texorpdfstring{$\widehat{B}_k = B_k$}{B\_k с шапкой = B\_k}: накопление}

Обновление отсчитывается от текущего приближения; вся информация,
усвоенная на предыдущих шагах, сохраняется в~$B_k$ неявно. Это
стандартный выбор для метода Бройдена и его вариантов; стоимость
шага~--- $O(n^2)$ на хранение и применение. В таком режиме локальная
q-сверхлинейная сходимость гарантируется условием Денниса--Море~\cite{dennis1974}~---
$\lim_{k\to\infty}\|(B_k - J(x^\ast))s_k\|/\|s_k\| = 0$,~--- а ошибка
$\|B_k - J(x^\ast)\|$ контролируется неравенством bounded
deterioration~\cite{broyden1973,dennis1996}; подробный анализ для
расширенного окна секущих см.\ в~гл.~\ref{ch:sp-broyden-theory}.

\subsection{\texorpdfstring{$\widehat{B}_k = B_0$}{B\_k с шапкой = B\_0}: пересборка}

Обновление отсчитывается от начальной матрицы (как правило, $B_0 = I$).
Информация за пределами текущего окна $m$~--- \textbf{стирается}:
$B_{k+1}$ зависит только от $B_0$ и последних $m$ пар $(s_j, y_j)$.

При $V_m = S_m$ и $B_0 = I$ формула~\eqref{eq:Vm_Sm} даёт:
\begin{equation}\label{eq:anderson}
  B_{k+1} = I + (Y_m - S_m)\,(S_m^\top S_m)^{-1}\,S_m^\top.
\end{equation}

\paragraph{Связь с ускорением Андерсона.}
\textbf{Ускорение Андерсона}~\cite{anderson1965}~--- классический метод
ускорения сходимости итераций с неподвижной точкой $x_{k+1} = G(x_k)$,
эквивалентных решению $F(x) = x - G(x) = 0$. На шаге $k$ оно строит новую
итерацию как линейную комбинацию последних $m+1$ точек: коэффициенты
$\alpha_i$ выбираются из условия $\sum_{i=0}^{m}\alpha_i = 1$ так,
чтобы минимизировать норму остатка
$\bigl\|\sum_{i=0}^{m}\alpha_i F(x_{k-i})\bigr\|$, после чего полученная
точка опционально смешивается с $G$ через параметр релаксации
$\beta\in(0,1]$ (при $\beta=1$~--- чистая интерполяционная схема).
Различают \emph{тип~I}, аппроксимирующий прямой якобиан, и
\emph{тип~II}~--- его обратный.

В~\cite{walker2011,fang2009} установлена эквивалентность: при
$B_0 = I$ и $\beta = 1$ мультисекущее обновление~\eqref{eq:anderson}
ранга~$m$ от единичной матрицы совпадает с шагом ускорения Андерсона
типа~I. Тем самым выбор $\widehat{B}_k = B_0$ вместо $\widehat{B}_k = B_k$
выводит из семейства методов Бройдена в семейство Андерсона~--- с
принципиально иным управлением исторической информацией: вместо
накопления через $B_k$ происходит полная пересборка от фиксированной базы
по скользящему окну из $m$ последних пар $(s_j, y_j)$.

\paragraph{Резюме.}
Общая мультисекантная формула~\eqref{eq:general} параметризуется двумя
независимыми объектами: матрицей направлений $V_m$, определяющей форму
обновления (минимум по норме Фробениуса при $V_m = S_m$, классический
метод Бройдена при $m=1$, $v_k=s_k$), и базовой матрицей $\widehat{B}_k$,
задающей режим работы с исторической информацией (накопление при
$\widehat{B}_k = B_k$ или пересборка при $\widehat{B}_k = B_0$,
эквивалентная ускорению Андерсона). В дальнейшем рассматривается режим
накопления с расширенным окном секущих условий~--- обновление SP-Broyden
строится в гл.~\ref{ch:sp-broyden-theory}.


%% ============================================================
